% =============================================================
%  CONCLUSIONES
% =============================================================
\renewcommand{\seccorta}{Conclusiones}
\section{Conclusiones}

\begin{frame}{Conclusiones}
  \begin{columns}[T]
    \column{0.48\textwidth}
    \begin{block}{Enlace A}
      \small
      \begin{itemize}
        \item \completar
        \item \completar
      \end{itemize}
    \end{block}
    \column{0.48\textwidth}
    \begin{block}{Enlace B}
      \small
      \begin{itemize}
        \item En zona urbana el obstáculo crítico son los edificios: la validación 3D en
              Google Earth fue indispensable porque UISP los subestima.
        \item Con \EBunoAlt/\EBdosAlt\ se despeja el \EBdespejePct\ de Fresnel ($A_{obs}=0$).
        \item $P_{RX}$ teórica \EBprxTeo\ vs.\ simulada \EBprxSim: diferencia $\approx$1.2 dB.
        \item SNR de \EBsnr\ $\Rightarrow$ 16QAM; capacidad $\approx$80 Mbps $>$ 50 Mbps requeridos.
      \end{itemize}
    \end{block}
  \end{columns}
  \vspace{4pt}
  \begin{alertblock}{Generales}
    \small El diseño es un compromiso entre capacidad, robustez frente al ruido y costo;
    la regulación (PIRE) acota la potencia utilizable. \completar[agregar en conjunto]
  \end{alertblock}
\end{frame}

\begin{frame}[plain]
  \centering\vfill
  {\usebeamercolor[fg]{title}\Huge\bfseries ¡Gracias!}\\[8pt]
  {\color{naranja}\large ¿Preguntas?}
  \vfill
\end{frame}
